% Paper B appendix, ported 2026-08-05 (merged-paper debt). Booktabs/width
% repairs applied; content otherwise verbatim from paperB_appendix.tex.
\section{Label definitions, nesting, and identifiability}\label{supporting-tables}

Every number in these tables is quoted in the body section that references it. They are placed here because the body has an eight-page limit and these tables document the results rather than making the argument.

\subsection{Observation modes}\label{observation-modes}

The grid \S\ref{sec:setup} refers to. The last column is the modality split Appendix~\ref{screenshot-tier} uses. It is not a cost split, for the reason \S\ref{sec:setup} gives: the image-bearing tier holds both the dearest mode and the cheapest.

\begin{table}[t]
\centering
\caption{The six observation modes (\S\ref{sec:setup}). The three P-modes keep either the mark legend, the SoM prompt, or both, while removing the per-step screenshot.}
\small
\setlength{\tabcolsep}{4pt}
\begin{tabular}{@{}llll@{}}
\toprule
mode & text payload & prompt family & annotated screenshot \\
\midrule
DOM & accessibility tree & DOM & no \\
SoM & mark legend & SoM & \textbf{yes} \\
Vision & none & vision & \textbf{yes} (unannotated) \\
P-text & mark legend & DOM & no \\
P-prompt & accessibility tree & SoM & no \\
P-SoM & mark legend & SoM & no \\
\bottomrule
\end{tabular}
\end{table}

\subsection{Where the which-mode label disagrees with measured cost}\label{where-the-which-mode-label-disagrees-with-measured-cost}

Per-cell detail for the 12.5--54.6\% range quoted in \S\ref{sec:lowerbound}, and for the claim that the exact-tie case never occurs.

\begin{table}[t]
\centering
\caption{Per-cell detail behind \S\ref{sec:lowerbound}. The last column counts labelled tasks on which the fixed priority list selected a mode strictly more expensive than another mode that also succeeded on that task. The exact-tie case, where the list order is the only tiebreaker, occurs in zero rows in every cell.}
\small
\setlength{\tabcolsep}{4pt}
\begin{tabular}{@{}llll@{}}
\toprule
cell & labels & multi-success & list picked a strictly pricier mode \\
\midrule
classifieds · B0 & 97 & 68 (70.1\%) & \textbf{53 (54.6\%)} \\
reddit · B0 & 53 & 36 (67.9\%) & \textbf{23 (43.4\%)} \\
classifieds · B1 & 55 & 29 (52.7\%) & \textbf{26 (47.3\%)} \\
reddit · B1 & 24 & 17 (70.8\%) & \textbf{9 (37.5\%)} \\
classifieds · B2 & 16 & 4 (25.0\%) & \textbf{2 (12.5\%)} \\
reddit · B2 & 15 & 3 (20.0\%) & \textbf{2 (13.3\%)} \\
\bottomrule
\end{tabular}
\end{table}

\subsection{Effect of nesting the operating-point selection}\label{effect-of-nesting-the-operating-point-selection}

Per-cell detail for the −0.99pp to +1.34pp range quoted in \S\ref{sec:lowerbound}, and for the observation that nesting moves the result in both directions.

\begin{table}[t]
\centering
\caption{Per-cell detail behind \S\ref{sec:lowerbound}. The naive column selects the threshold and the strongest and cheapest modes from whole-cell outcomes; the nested column re-derives all three inside each outer fold. Nesting moves the result in both directions.}
\small
\setlength{\tabcolsep}{4pt}
\begin{tabular}{@{}llll@{}}
\toprule
cell & naive nesting SR & fully nested SR & Δ \\
\midrule
classifieds · B0 & 25.45\% & 26.79\% & \textbf{+1.34pp} \\
classifieds · B1 & 13.84\% & 14.29\% & +0.45pp \\
classifieds · B2 & 1.34\% & 1.34\% & ±0.00pp \\
reddit · B0 & 13.79\% & 12.81\% & \textbf{−0.99pp} \\
reddit · B1 & 6.40\% & 5.91\% & −0.49pp \\
reddit · B2 & 3.94\% & 3.94\% & ±0.00pp \\
\bottomrule
\end{tabular}
\end{table}

\subsection{Identifiability of pooled which-mode labels}\label{identifiability-of-pooled-which-mode-labels}

Per-site detail for the conflict rates and modal-agreement figures quoted in Appendix~\ref{pooling-across-backbones}.

\begin{table}[t]
\centering
\caption{Per-site detail behind Appendix~\ref{pooling-across-backbones}. A conflict is one task on which two cells recorded different oracle modes. In-sample modal agreement is the accuracy of emitting the modal label per task, over all pooled labelled rows (168 on classifieds, 92 on reddit), a row being one (task, backbone) pair on which that backbone solved the task.}
\small
\setlength{\tabcolsep}{4pt}
\begin{tabular}{@{}
  >{\raggedright\arraybackslash}p{(\linewidth - 12\tabcolsep) * \real{0.1429}}
  >{\raggedright\arraybackslash}p{(\linewidth - 12\tabcolsep) * \real{0.1429}}
  >{\raggedright\arraybackslash}p{(\linewidth - 12\tabcolsep) * \real{0.1429}}
  >{\raggedright\arraybackslash}p{(\linewidth - 12\tabcolsep) * \real{0.1429}}
  >{\raggedright\arraybackslash}p{(\linewidth - 12\tabcolsep) * \real{0.1429}}
  >{\raggedright\arraybackslash}p{(\linewidth - 12\tabcolsep) * \real{0.1429}}
  >{\raggedright\arraybackslash}p{(\linewidth - 12\tabcolsep) * \real{0.1429}}@{}}
\toprule
\begin{minipage}[b]{\linewidth}\raggedright
site
\end{minipage} & \begin{minipage}[b]{\linewidth}\raggedright
tasks labelled in ≥2 cells
\end{minipage} & \begin{minipage}[b]{\linewidth}\raggedright
conflicting
\end{minipage} & \begin{minipage}[b]{\linewidth}\raggedright
conflict rate
\end{minipage} & \begin{minipage}[b]{\linewidth}\raggedright
in-sample modal agreement
\end{minipage} & \begin{minipage}[b]{\linewidth}\raggedright
same, on shared tasks only
\end{minipage} & \begin{minipage}[b]{\linewidth}\raggedright
same, exact-vector grouping
\end{minipage} \\
\midrule
classifieds & 54 & 31 & \textbf{57.4\%} & \textbf{79.2\%} & 70.3\% (n=118) & 83.9\% \\
reddit & 25 & 14 & \textbf{56.0\%} & \textbf{83.7\%} & 74.1\% (n=58) & 89.1\% \\
\bottomrule
\end{tabular}
\end{table}

\emph{We call it agreement and not a Bayes ceiling because it is a resubstitution estimate: it scores the same rows it took the modal label from, so a task labelled by only one backbone is correct by construction. That describes 50 of 168 classifieds rows (29.8\%) and 34 of 92 reddit rows (37.0\%); restricted to tasks two or more backbones label, agreement falls to 70.3\% and 74.1\%. An out-of-sample bound would need leave-one-backbone-out prediction or a shrinkage estimator, and would be lower still. Every number here is therefore an optimistic bound on what a pooled classifier could reach, which is the direction the argument needs.}

\emph{The last column groups by the exact feature vector instead of by task. Rows of one task are not always identical, because three of the five numeric features are read from that backbone's own step-0 observation, so they differ somewhere on 31.5\% of shared classifieds tasks and 80.0\% of shared reddit tasks. We report but do not use that grouping: it leaves 74 of 117 classifieds groups and 69 of 78 reddit groups with a single member, covering 44\% and 75\% of rows, so it is even more inflated than the headline, and a router serving one backbone could not recover backbone identity from that jitter anyway.}

\subsection{The screenshot-modality tier}\label{the-screenshot-modality-tier}

Per-site detail for Appendix~\ref{screenshot-tier}, on the same pooled labelled rows and the same grouping as Appendix~\ref{identifiability-of-pooled-which-mode-labels}. The agreement column is the exception: it is defined only on tasks labelled by two or more backbones, since agreement needs two labels to compare.

\begin{table}[t]
\centering
\caption{Per-site detail behind Appendix~\ref{screenshot-tier}, over the same solve events. Columns two and three carry the same resubstitution caveat as Appendix~\ref{identifiability-of-pooled-which-mode-labels}, and column three additionally rises for an arithmetic reason: merging six classes into two can only increase a modal share. The claim that backbones agree about the screenshot therefore rests on the last two columns, which measure agreement between two backbones' labels directly rather than against a modal label. Under Appendix~\ref{identifiability-of-pooled-which-mode-labels}'s exact-vector grouping the tier figures are 92.3\% and 97.8\%. No classifier is fitted to this target anywhere in the paper.}
\small
\setlength{\tabcolsep}{4pt}
\begin{tabular}{@{}
  >{\raggedright\arraybackslash}p{(\linewidth - 8\tabcolsep) * \real{0.2000}}
  >{\raggedright\arraybackslash}p{(\linewidth - 8\tabcolsep) * \real{0.2000}}
  >{\raggedright\arraybackslash}p{(\linewidth - 8\tabcolsep) * \real{0.2000}}
  >{\raggedright\arraybackslash}p{(\linewidth - 8\tabcolsep) * \real{0.2000}}
  >{\raggedright\arraybackslash}p{(\linewidth - 8\tabcolsep) * \real{0.2000}}@{}}
\toprule
\begin{minipage}[b]{\linewidth}\raggedright
site
\end{minipage} & \begin{minipage}[b]{\linewidth}\raggedright
which-mode modal agreement
\end{minipage} & \begin{minipage}[b]{\linewidth}\raggedright
tier modal agreement
\end{minipage} & \begin{minipage}[b]{\linewidth}\raggedright
tier agreement across backbones
\end{minipage} & \begin{minipage}[b]{\linewidth}\raggedright
six-way agreement across backbones
\end{minipage} \\
\midrule
classifieds & 79.2\% & \textbf{89.9\%} & \textbf{68.5\%} & 42.6\% \\
reddit & 83.7\% & \textbf{96.7\%} & \textbf{88.0\%} & 44.0\% \\
\bottomrule
\end{tabular}
\end{table}

\section{Derivations for the four relabelling routes}\label{derivations-for-the-four-relabelling-routes}

\S\ref{sec:lowerbound} and \S\ref{sec:gap} state the outcome of each route. This appendix gives the derivation.

\subsection{Continuous labels}\label{continuous-labels}

The cleanest fix would be to regress on a graded quality signal rather than classify a discrete winner: partial credit turns every episode into a training example regardless of whether it succeeded. VisualWebArena does not provide one. Across the 7,686 scored episodes of our 36 landed conditions the evaluator emits exactly two values, 0.0 and 1.0, in a 7,041 / 645 split, as do the 36 protocol-excluded episodes we do not score. This is a property of the benchmark's evaluation design rather than of our pipeline, and it forecloses the route entirely.

\subsection{Coarser classes}\label{coarser-classes}

\S\ref{sec:lowerbound}'s obstruction is that four of six cells have fewer than ten labelled rows in more than one class. Merging classes does not add rows. The binary collapse of the screenshot tier (Appendix~\ref{screenshot-tier}) does help, but for a different reason, having to do with agreement across backbones rather than with class count.

\subsection{Pooling across backbones}\label{pooling-across-backbones}

Six cells at 15--97 labels become 260 pooled examples and every class clears the minimum-count filter, so supply is solved. Identifiability is not. The features carry no model identity, so two backbones facing the same task on the same site produce near-identical feature vectors, and where they are identical and the oracle labels differ, a classifier is being asked to emit two different answers for one input. The figure reported in Appendix~\ref{identifiability-of-pooled-which-mode-labels} is the accuracy of emitting the modal label per group, scored on the rows the label came from. It is an in-sample bound on what a pooled classifier could reach, not a Bayes ceiling, and Appendix~\ref{identifiability-of-pooled-which-mode-labels} gives the resubstitution caveat.

They are only near-identical, not identical, because \texttt{dom\_\allowbreak{}complexity}, \texttt{text\_\allowbreak{}length} and \texttt{tokens\_\allowbreak{}input\_\allowbreak{}text} are read from the backbone's own step-0 observation rather than from the task config. On 31.5\% of shared classifieds tasks and 80.0\% of shared reddit tasks the rows therefore differ somewhere. Grouping by the exact vector rather than by task raises the figure (Appendix~\ref{identifiability-of-pooled-which-mode-labels}, last column), and we report but do not adopt that number: it leaves most groups with one member, and a group of one is scored perfectly whatever the labels do, so it inflates with feature sparsity rather than tracking identifiability. Every version of the number is an optimistic bound and all of them are far below what a deployable which-mode router needs, which is the only thing the argument turns on.

\subsection{Screenshot tier}\label{screenshot-tier}

The tier label is derived from the same solve events as the which-mode label, by mapping each oracle mode to image-bearing (SoM, Vision) or text-only (DOM, P-text, P-prompt, P-SoM). No new episodes are involved, which is why the figure rises without a single new solve event. Part of that rise is arithmetic (two classes admit a larger modal share than six), so the agreement columns rather than the modal-agreement columns carry the claim. The tier is defined only on tasks that some mode solved, so its denominator is the solved set and not the full task universe. Appendix~\ref{the-screenshot-modality-tier}'s modal-agreement columns run over the whole pooled labelled set, as in Appendix~\ref{identifiability-of-pooled-which-mode-labels}; the two agreement columns are restricted to tasks labelled by two or more backbones, that being the only set on which cross-backbone agreement is defined.

\section{The reddit · B2 saving in detail}\label{the-reddit-b2-saving-in-detail}

\S\ref{sec:lowerbound} reports that the one cell whose cost saving survives Holm correction has an AUROC below chance. The mechanism is tail enrichment rather than a globally ordered score.

Reddit · B2 sends 192 of 203 tasks (95\%) to the cheap mode with no accuracy loss, which is unsurprising in a cell where only 7.4\% of tasks are solvable at all: almost nothing in that 95\% was going to succeed under any mode. The policy therefore differs from the free always-cheapest policy by five percent of the allocation. The 11 tasks it holds back for the strong mode carry four successes that the fixed policy does not collect, against four collected by the fixed policy overall. The permutation null detects that enrichment. A globally ordered score is not required to produce it, which is why the cell's AUROC of 0.483, below both chance and its own best single covariate at 0.711, is consistent with a real saving.

Two properties of that test are worth recording. It runs 10,000 draws and reports the plus-one Monte Carlo estimator (k+1)/(B+1), whose floor is therefore 1.0 × 10⁻⁴, two orders below the tightest Holm threshold of 8.3 × 10⁻³. That matters because at the 200 draws we first used, the floor was 1/201 = 4.98 × 10⁻³, this cell reported exactly it, and whether it could clear the threshold at all was a function of the draw count rather than of the data; at 10,000 draws four of the draws match or beat the observed saving, so p = 5.0 × 10⁻⁴ is measured. Second, the quantity tested is the saving at an operating point selected against whole-cell outcomes, which is not the nested policy of Table 6. Null and observation select that point the same way, so the null absorbs the selection optimism and the comparison is fair, but the point is not one a deployment could occupy, and \S\ref{sec:lowerbound}'s conclusion rests on the nested numbers rather than on this test.

\subsection{Supply and trainability under both label definitions}\label{supply-and-trainability-under-both-label-definitions}

\S\ref{sec:lowerbound} keeps the prior-order label and reports the measured-cost rule as a sensitivity. Supply is identical under both by construction, since each labels exactly the tasks some mode solved, so only the class distribution and through it trainability can move. The relabelled column equals the ``order picked a strictly pricier mode'' column of A.2 exactly, which is the consistency check: a label moves if and only if the order's pick was not the measured cheapest.

\begin{table}[t]
\centering
\caption{Trainability under the two label definitions. ``Surviving'' counts classes clearing ten training rows in a five-fold split; \textbf{no} marks a cell with fewer than two, where no classifier exists. Four of six cells are untrainable under the reported label and five of six under the measured-cost alternative, so the supply argument does not depend on the choice.}
\small
\setlength{\tabcolsep}{4pt}
\begin{tabular}{@{}
  >{\raggedright\arraybackslash}p{(\linewidth - 8\tabcolsep) * \real{0.2000}}
  >{\raggedright\arraybackslash}p{(\linewidth - 8\tabcolsep) * \real{0.2000}}
  >{\raggedright\arraybackslash}p{(\linewidth - 8\tabcolsep) * \real{0.2000}}
  >{\raggedright\arraybackslash}p{(\linewidth - 8\tabcolsep) * \real{0.2000}}
  >{\raggedright\arraybackslash}p{(\linewidth - 8\tabcolsep) * \real{0.2000}}@{}}
\toprule
\begin{minipage}[b]{\linewidth}\raggedright
cell
\end{minipage} & \begin{minipage}[b]{\linewidth}\raggedright
labels
\end{minipage} & \begin{minipage}[b]{\linewidth}\raggedright
relabelled
\end{minipage} & \begin{minipage}[b]{\linewidth}\raggedright
prior order: surviving classes
\end{minipage} & \begin{minipage}[b]{\linewidth}\raggedright
measured cost: surviving classes
\end{minipage} \\
\midrule
classifieds · B0 & 97 & 53 & 3 (DOM, P-prompt, SoM) & 2 (SoM, Vision) \\
reddit · B0 & 53 & 23 & 1 (DOM) (\textbf{no}) & 1 (P-text) (\textbf{no}) \\
classifieds · B1 & 55 & 26 & 2 (DOM, SoM) & 1 (Vision) (\textbf{no}) \\
reddit · B1 & 24 & 9 & 0 (\textbf{no}) & 0 (\textbf{no}) \\
classifieds · B2 & 16 & 2 & 0 (\textbf{no}) & 0 (\textbf{no}) \\
reddit · B2 & 15 & 2 & 0 (\textbf{no}) & 0 (\textbf{no}) \\
\bottomrule
\end{tabular}
\end{table}

The single cell that changes, classifieds · B1, loses a class rather than gaining one: the prior-order label keeps DOM and SoM above the threshold, the measured-cost label concentrates enough of those rows onto Vision that only Vision survives. The alternative definition therefore strengthens the negative result, which is a reason to report it and not a reason to adopt it.

\subsection{The best-success mode is not stable across folds}\label{the-best-success-mode-is-not-stable-across-folds}

The nested design of \S\ref{sec:lowerbound} re-selects the best-success mode inside every outer fold, which exposes something the whole-cell version conceals. In reddit · B0 the five outer folds select DOM, DOM, SoM, SoM, DOM. A pipeline that picks one best mode from all realised outcomes is therefore not merely optimistic about its threshold; it reports a mode choice that its own resampling does not reproduce.
